\section{Motivation}

If vectors are the standard container for data, \emph{matrices} are the standard container for everything
that is two-dimensional — or that \emph{acts} on data. This makes them, with little exaggeration, the
central data structure of computing:

\begin{itemize}
\item a greyscale \textbf{image} is a matrix of pixel intensities, and image filters (blur, sharpen, edge
detection) are matrix operations;
\item a \textbf{graph} — a road network, a social network, the web — is stored as its \emph{adjacency
matrix}, whose $ij$-entry records the connection between node $i$ and node $j$; counting walks, routing
flows and ranking nodes all become matrix computations (Chapter~\ref{ch:eigen} makes this concrete with
Google's PageRank);
\item a layer of a \textbf{neural network} is a matrix of weights: an input vector is multiplied by the
matrix, a simple non-linear function is applied, and the result is passed to the next layer.
``Training'' a network means nothing else than adjusting the entries of its matrices;
\item a \textbf{spreadsheet}, a database table, the term--document table inside a search engine — matrices
again.
\end{itemize}

The strangest-looking definition of this chapter — the row-by-column rule for \emph{multiplying} matrices —
is also its most consequential one. It is designed precisely so that multiplication of matrices corresponds
to \emph{composition} of transformations (we prove this in Chapter~\ref{ch:linear-mappings}), and its
central role in graphics and machine learning explains a striking fact about modern hardware: GPUs, the
chips that power both video games and the training of large AI models, are at heart machines optimised for
one task — multiplying enormous matrices fast. Even the theory is still alive: the true asymptotic
complexity of matrix multiplication (naively $O(n^3)$, improved to about $O(n^{2.81})$ by Strassen's
recursive algorithm) remains a famous open problem of theoretical computer science.

Towards the end of the chapter we develop \emph{elementary operations}, \emph{echelon forms} and the
notion of \emph{rank}. Rank measures the genuine information content of a matrix: a $1000\times 1000$
matrix of rank $20$ can be reconstructed from roughly $40\,000$ numbers instead of a million. Low rank is
the compressible structure hiding in most real-world data — and the algorithm that computes it, row
reduction, is the same engine that will solve systems of equations in Chapter~\ref{ch:sole} and invert
matrices later in this chapter.

The chapter closes with the \emph{determinant}, which compresses an entire $n\times n$ matrix into a single scalar that answers a remarkable
number of yes/no questions at once. Is the matrix invertible? Does the system $AX=B$ have exactly one
solution? Do these $n$ vectors form a basis of $\K^n$? Does the transformation collapse space onto
something flatter? All of these questions are equivalent, and all reduce to checking $\det(A)\neq 0$ —
which is why a determinant test appears in virtually every algorithm that must know whether a matrix is
``degenerate''.

For a computer scientist the determinant has a second life as a \emph{geometric} quantity:

\begin{itemize}
\item In 2D, $|\det(A)|$ is the \textbf{area} of the parallelogram spanned by the rows of $A$; in 3D it is
a volume. A transformation with determinant $6$ magnifies every area six-fold; determinant $0$ means space
gets flattened.
\item The \textbf{sign} of the determinant detects \emph{orientation}. The test ``does point $C$ lie to
the left or to the right of the line through $A$ and $B$?'' is the sign of a $2\times 2$ determinant; it
is the basic primitive of computational geometry, executed millions of times inside convex-hull
algorithms, triangulations, collision detection and map matching. In 3D rendering the same sign test
(\emph{backface culling}) lets the GPU skip the triangles that face away from the camera.
\item In graphics and machine learning, determinants of \emph{Jacobian} matrices track how transformations
distort volume — the bookkeeping device behind texture mapping and behind generative ``normalizing flow''
models.
\end{itemize}

Part of the motivation is also a computational warning. The Laplace expansion defined below is
mathematically fundamental but computationally explosive: applied naively it performs on the order of $n!$
operations. The behaviour of determinants under row operations, established in this chapter, brings the
cost down to that of Gaussian elimination, $O(n^3)$ — one of the nicest examples in this course of theory
rescuing an algorithm. The same warning applies to the elegant closed formulas (the adjugate formula
for $A^{-1}$ in this chapter, and Cramer's rule in the next): they are invaluable for insight and perfectly practical for small
matrices, but for large ones the method of choice is, once again, row reduction.

\section{Matrices}

\begin{definition}
A \emph{matrix of size $m\times n$ over a field $\K$} (or simply a \emph{matrix}) is a rectangular table of
elements $a_{ij}\in\K$ for $i=1,\ldots,m$, $j=1,\ldots,n$:
\[
\left( \begin{array}{cccc}
a_{11} & a_{12} & \ldots & a_{1n} \\
a_{21} & a_{22} & \ldots & a_{2n} \\
\vdots & \vdots & \ddots & \vdots \\
a_{m1} & a_{m2} & \ldots & a_{mn}
\end{array} \right).
\]
We denote this matrix by $(a_{ij})_{m\times n}$ or simply $(a_{ij})$. The element $a_{ij}$ is called the
\emph{$ij$-entry}. Matrices are usually denoted by capital letters $A,B,\ldots$. Two matrices $A,B$ are
\emph{equal} (written $A=B$) if they have the same size and the same $ij$-entries. The set of matrices of size
$m\times n$ over $\K$ is denoted by $\K_m^n$.
\end{definition}

\begin{definition}
Given a matrix $(a_{ij})_{m\times n}$, the $n$-tuples
\[
(a_{11}, a_{12}, \ldots, a_{1n}),\ \ldots,\ (a_{m1}, a_{m2}, \ldots, a_{mn})
\]
are its \emph{rows}, and the $m$-tuples
\[
\left( \begin{array}{c} a_{11}\\ a_{21} \\ \vdots \\ a_{m1} \end{array} \right), \ldots,
\left( \begin{array}{c} a_{1n} \\ a_{2n} \\ \vdots \\ a_{mn} \end{array} \right)
\]
are its \emph{columns}.
\end{definition}

\begin{example}
The matrix
\[
\left( \begin{array}{ccc}
1 & 0 & 2 \\
-3 & \tfrac{1}{2} & 4
\end{array} \right)
\]
is a $2\times 3$ matrix over $\R$. Its rows are $(1 , 0 , 2)$ and $(-3 , \tfrac{1}{2}, 4)$, and its columns are
\[
\left( \begin{array}{c} 1 \\ -3 \end{array} \right),\quad
\left( \begin{array}{c} 0 \\ \tfrac{1}{2} \end{array} \right),\quad
\left( \begin{array}{c} 2 \\ 4 \end{array} \right).
\]
\end{example}

\begin{fact}
Let $A$ be an $m\times n$ matrix over a field $\K$. The rows of $A$ are members of the vector space $\K^n$ and
the columns are members of $\K^m$.
\end{fact}

\begin{definition}
An $m\times n$ matrix is called a \emph{square matrix} if $m=n$. A square matrix is \emph{symmetric} if
$a_{ij}=a_{ji}$ for all $i,j=1,\ldots,n$. The $n$-tuple $(a_{11},a_{22},\ldots,a_{nn})$ is called the
\emph{main diagonal} of a square matrix $(a_{ij})_{n\times n}$.
\end{definition}

\begin{example}
The following are square matrices, the second one symmetric:
\[
\left( \begin{array}{cc} 1 & 0 \\ -3 & \tfrac{1}{2} \end{array} \right), \qquad
\left( \begin{array}{cc} 1 & -3 \\ -3 & 1 \end{array} \right).
\]
\end{example}

\begin{definition}
A square matrix $A = (a_{ij})_{n\times n}$ is called
\begin{itemize}
\item \emph{diagonal} if $a_{ij} = 0$ for all $i\neq j$;
\item \emph{upper triangular} if $a_{ij} = 0$ for all $i>j$ (all entries below the main diagonal are zero);
\item \emph{lower triangular} if $a_{ij} = 0$ for all $i<j$.
\end{itemize}
\end{definition}

\begin{example}
The matrices
\[
\left( \begin{array}{cc} 2 & 0 \\ 0 & -1 \end{array} \right), \qquad
\left( \begin{array}{ccc} 1 & 5 & 0 \\ 0 & 2 & -1 \\ 0 & 0 & 3 \end{array} \right)
\]
are diagonal and upper triangular, respectively. Every diagonal matrix is both upper and lower triangular.
\end{example}

\section{Matrix operations}

\begin{definition}[Addition]
Let $A=(a_{ij})$ and $B=(b_{ij})$ be two matrices of the same size $m\times n$. The \emph{sum} of $A$ and $B$,
denoted $A+B$, is the matrix
\[
A+B = (a_{ij}+b_{ij})_{m\times n}.
\]
\end{definition}

\begin{definition}[Scalar multiplication]
The \emph{product} of a matrix $A=(a_{ij})$ by a scalar $k\in \K$, denoted $k\cdot A$, is the matrix
$k\cdot A = (k\,a_{ij})$. We write $-A := (-1)\cdot A$ and $A-B := A+(-B)$.
\end{definition}

\begin{example}
\[
\left( \begin{array}{ccc} 1 & 0 & 2 \\ -3 & \tfrac{1}{2} & 4 \end{array} \right)
+
\left( \begin{array}{ccc} -3 & 1 & 4 \\ 2 & \tfrac{3}{2} & 0 \end{array} \right) =
\left( \begin{array}{ccc} -2 & 1 & 6 \\ -1 & 2 & 4 \end{array} \right),
\qquad
3\cdot \left( \begin{array}{ccc} 1 & 0 & 2 \\ -3 & \tfrac{1}{2} & 4 \end{array} \right)
=
\left( \begin{array}{ccc} 3 & 0 & 6 \\ -9 & \tfrac{3}{2} & 12 \end{array} \right).
\]
\end{example}

\begin{warning}
Addition of matrices of different sizes is \emph{not} defined.
\end{warning}

\begin{definition}
The $m\times n$ matrix with all entries equal to zero is called the \emph{zero matrix} and is denoted by
$0_{m\times n}$, or simply $0$ when the size is clear from the context.
\end{definition}

\begin{theorem}
Let $A,B,C$ be $m\times n$ matrices over a field $\K$ and let $k_1,k_2\in \K$. Then
\[
\begin{aligned}
&A+0 = 0+A =A, &&k_1(A+B) = k_1\cdot A+k_1\cdot B,\\
&(A+B)+C = A+(B+C), &&(k_1+k_2)A = k_1\cdot A+k_2\cdot A,\\
&A+B = B+A, &&(k_1 k_2)A = k_1(k_2\cdot A),\\
&A-A = 0, && 1\cdot A = A.
\end{aligned}
\]
\end{theorem}

\begin{definition}[Matrix multiplication]
Let $A$ be a $k\times m$ matrix and $B$ an $m\times n$ matrix. The \emph{product} $A\cdot B$ is the
$k\times n$ matrix $C=(c_{ij})$ whose entries are
\[
c_{ij} = a_{i1}b_{1j}+a_{i2}b_{2j}+\ldots + a_{im}b_{mj}, \qquad i=1,\ldots,k,\ \ j=1,\ldots,n.
\]
\end{definition}

\begin{warning}
The product $A\cdot B$ is defined \emph{only} when the number of columns of $A$ equals the number of rows of
$B$.
\end{warning}

\begin{example}
Let
\[
A = \left( \begin{array}{ccc} 1 & 0 & 2 \\ 2 & 1 & 3 \end{array} \right), \qquad
B = \left( \begin{array}{cc} 1 & 2 \\ 0 & 1 \\ 1 & 0 \end{array} \right).
\]
Here $A$ is $2\times 3$ and $B$ is $3\times 2$, so both products are defined: $A\cdot B$ is $2\times 2$ and
$B\cdot A$ is $3\times 3$. We compute
\[
A\cdot B = \left( \begin{array}{cc}
1\cdot 1+0\cdot 0+2\cdot 1 & 1\cdot 2+0\cdot 1+2\cdot 0 \\
2\cdot 1+1\cdot 0+3\cdot 1 & 2\cdot 2+1\cdot 1+3\cdot 0
\end{array} \right)
= \left( \begin{array}{cc} 3 & 2 \\ 5 & 5 \end{array} \right),
\qquad
B\cdot A = \left( \begin{array}{ccc} 5 & 2 & 8 \\ 2 & 1 & 3 \\ 1 & 0 & 2 \end{array} \right).
\]
In particular, $A\cdot B$ and $B\cdot A$ need not even have the same size.
\end{example}

\begin{remark}[Matrix--vector products]
An important special case is when the second factor is a single column: the product of an $m\times n$
matrix $A$ with a column vector $\vv\in\K^n$ is a column vector $A\vv\in\K^m$. Writing the definition out reveals a useful reading:
\[
A \left(\begin{array}{c} x_1\\ \vdots\\ x_n\end{array}\right)
= x_1\cdot \big(\text{1st column of } A\big) + \ldots + x_n\cdot \big(\text{$n$-th column of } A\big),
\]
i.e.\ $A\vv$ is the \emph{linear combination of the columns of $A$} with coefficients the entries of
$\vv$. This observation will be used repeatedly in Chapters~\ref{ch:sole} and~\ref{ch:linear-mappings}.
\end{remark}

\begin{remark}
All of matrix algebra works verbatim over \emph{any} field $\K$ — including the finite fields $\Z_p$ —
and the same is true of everything that follows in this book: row reduction, rank, systems of linear
equations, determinants and inverses. Over $\Z_p$ one simply reduces modulo $p$ at every step.
\end{remark}

\begin{example}[Matrices over a finite field]
Over the field $\Z_5$,
\[
\left( \begin{array}{cc} 2 & 3 \\ 4 & 1 \end{array} \right)
\left( \begin{array}{cc} 1 & 4 \\ 2 & 3 \end{array} \right)
= \left( \begin{array}{cc} 2+6 & 8+9 \\ 4+2 & 16+3 \end{array} \right)
= \left( \begin{array}{cc} 3 & 2 \\ 1 & 4 \end{array} \right),
\]
since $8 \equiv 3$, $17 \equiv 2$, $6\equiv 1$ and $19 \equiv 4$ modulo $5$.
\end{example}

\begin{definition}
The square matrix of size $n\times n$ of the form
\[
\left( \begin{array}{cccc}
1 & 0 & \ldots & 0 \\
0 & 1 & \ldots & 0 \\
\vdots & \vdots & \ddots & \vdots \\
0 & 0 & \ldots & 1
\end{array} \right)
\]
is called the \emph{identity matrix} and is denoted by $I_{n\times n}$ or $I$.
\end{definition}

\begin{theorem}[Properties of multiplication]
Let $A,A_1$ be $k\times m$ matrices, $B,B_1$ be $m\times n$ matrices and $C$ be an $n\times o$ matrix. Then
\[
\begin{aligned}
& (A\cdot B)\cdot C = A\cdot (B\cdot C),\\
& A\cdot I_{m\times m} = I_{k\times k} \cdot A = A,\\
& (A+A_1)\cdot B = A \cdot B +A_1\cdot B,\\
& A\cdot (B+B_1) = A \cdot B +A\cdot B_1.
\end{aligned}
\]
\end{theorem}

\begin{warning}
For $A,B\in \K_n^n$ multiplication is in general \emph{not commutative}: sometimes $A\cdot B \neq B\cdot A$.
\end{warning}

\begin{example}
Even for square matrices the order of multiplication matters:
\[
\left( \begin{array}{cc} 0 & 1 \\ 0 & 0 \end{array} \right)
\left( \begin{array}{cc} 0 & 0 \\ 1 & 0 \end{array} \right)
= \left( \begin{array}{cc} 1 & 0 \\ 0 & 0 \end{array} \right),
\qquad
\left( \begin{array}{cc} 0 & 0 \\ 1 & 0 \end{array} \right)
\left( \begin{array}{cc} 0 & 1 \\ 0 & 0 \end{array} \right)
= \left( \begin{array}{cc} 0 & 0 \\ 0 & 1 \end{array} \right).
\]
\end{example}

\begin{warning}
A product of two non-zero matrices can be the zero matrix: for
$A = \left( \begin{smallmatrix} 0 & 1 \\ 0 & 0 \end{smallmatrix} \right)$ we have $A\neq 0$ but
$A\cdot A = 0$. Consequently, one cannot cancel in matrix equations: $A\cdot B = A\cdot C$ does \emph{not}
imply $B = C$.
\end{warning}

\begin{definition}[Matrix powers]
For a square matrix $A$ we set $A^0 = I$ and $A^{k} = A^{k-1}\cdot A$ for $k\geq 1$. Since matrix
multiplication is associative, the usual rules $A^k\cdot A^l = A^{k+l}$ and $(A^k)^l = A^{kl}$ hold.
\end{definition}

\begin{example}[Counting walks in a graph]
Let $A$ be the \emph{adjacency matrix} of a graph: $a_{ij} = 1$ if there is an edge from node $i$ to node
$j$, and $a_{ij} = 0$ otherwise. Then the $ij$-entry of $A^k$ counts the walks of length $k$ from $i$ to
$j$. For $k = 2$ this is immediate from the definition of the product: the entry
$\sum_{m} a_{im}\,a_{mj}$ gains $1$ exactly for those intermediate nodes $m$ with a step $i\to m$ and a
step $m\to j$; the general case follows by induction. This is the walk-counting claim of the motivation
section made precise.
\end{example}

\begin{definition}[Transpose]
Let $A=(a_{ij})$ be an $m\times n$ matrix. The \emph{transpose} of $A$, denoted $A^T$, is the $n\times m$
matrix obtained by interchanging rows and columns:
\[
\left( \begin{array}{cccc}
a_{11} & a_{12} & \ldots & a_{1n} \\
a_{21} & a_{22} & \ldots & a_{2n} \\
\vdots & \vdots & \ddots & \vdots \\
a_{m1} & a_{m2} & \ldots & a_{mn}
\end{array} \right)^T =
\left( \begin{array}{cccc}
a_{11} & a_{21} & \ldots & a_{m1} \\
a_{12} & a_{22} & \ldots & a_{m2} \\
\vdots & \vdots & \ddots & \vdots \\
a_{1n} & a_{2n} & \ldots & a_{mn}
\end{array} \right).
\]
\end{definition}

\begin{example}
\[
\left( \begin{array}{ccc} 1 & 0 & 2 \\ -3 & \tfrac{1}{2} & 4 \end{array} \right)^T =
\left( \begin{array}{cc} 1 & -3 \\ 0 & \tfrac{1}{2} \\ 2 & 4 \end{array} \right).
\]
\end{example}

\begin{fact}
A square matrix $A$ is symmetric if and only if $A^T=A$.
\end{fact}

\begin{theorem}[Properties of the transpose]
Let $A,B$ be matrices of appropriate sizes and $k$ a scalar. Then
\[
(A+B)^T = A^T + B^T, \qquad (A^T)^T = A, \qquad (k\cdot A)^T = k\cdot A^T, \qquad (A\cdot B)^T = B^T\cdot A^T.
\]
\end{theorem}

\section{Elementary operations}

\begin{definition}
By \emph{elementary row operations} on a matrix $A$ we mean the following:
\begin{enumerate}
\item \textbf{(Row switching)} the $i$-th and $j$-th rows are interchanged ($R_i\leftrightarrow R_j$);
\item \textbf{(Row scaling)} each element of the $i$-th row is multiplied by a \emph{nonzero} scalar
$k\in \K$ ($kR_i\rightarrow R_i$);
\item \textbf{(Row addition)} the $i$-th row is replaced by the sum of the $i$-th row and a multiple of the
$j$-th row ($R_i+k\cdot R_j\rightarrow R_i$).
\end{enumerate}
\emph{Elementary column operations} are defined analogously.
\end{definition}

\begin{example}
\[
\left( \begin{array}{cccc}
-3 & 1 & 4 & 2\\ 2 & \tfrac{3}{2} & 0 & 3\\ -3 & 1 & 2 & 1
\end{array} \right)
\xrightarrow{R_1\leftrightarrow R_2}
\left( \begin{array}{cccc}
2 & \tfrac{3}{2} & 0 & 3\\ -3 & 1 & 4 & 2\\ -3 & 1 & 2 & 1
\end{array} \right)
\xrightarrow{3R_2\to R_2}
\left( \begin{array}{cccc}
2 & \tfrac{3}{2} & 0 & 3\\ -9 & 3 & 12 & 6\\ -3 & 1 & 2 & 1
\end{array} \right)
\xrightarrow{R_2+2R_1\to R_2}
\left( \begin{array}{cccc}
2 & \tfrac{3}{2} & 0 & 3\\ -5 & 6 & 12 & 12\\ -3 & 1 & 2 & 1
\end{array} \right).
\]
\end{example}

\begin{theorem}
Each elementary row (column) operation is invertible, and its inverse is again an elementary row (resp.\
column) operation.
\end{theorem}

\begin{proof}
\begin{enumerate}
\item The inverse of $R_i\leftrightarrow R_j$ is $R_j\leftrightarrow R_i$.
\item The inverse of $k\cdot R_i\rightarrow R_i$ (with $k\neq 0$) is $\tfrac{1}{k}\cdot R_i\rightarrow R_i$.
\item The inverse of $R_i+k\cdot R_j\rightarrow R_i$ is $R_i-k\cdot R_j\rightarrow R_i$.
\end{enumerate}
Replacing the letter $R$ by $C$ gives the proof for elementary column operations.
\end{proof}

\begin{definition}
Two matrices $A$ and $B$ of the same size are \emph{row equivalent} if one can be obtained from the other
using elementary row operations.
\end{definition}

\section{Row echelon form and rank}

\begin{definition}
A matrix $A$ is in \emph{row echelon form} if
\begin{enumerate}
\item all zero rows are at the bottom, and
\item the first nonzero entry in the $i$-th row — called the \emph{leading entry} (or \emph{pivot}) of that
row — lies to the right of the leading entry of the row above it.
\end{enumerate}
\end{definition}

\begin{example}
The matrix
\[
\left( \begin{array}{cccc}
2 & 1 & 0 & 3\\ 0 & 6 & 12 & 12\\ 0 & 0 & 2 & 1 \\ 0 & 0 & 0 & 0
\end{array} \right)
\]
is in row echelon form.
\end{example}

\begin{fact}
Any matrix $A$ is row equivalent to a matrix in row echelon form.
\end{fact}

\begin{example}
\[
\left( \begin{array}{ccc}
1 & 0 & 2\\ 2 & -1 & 3\\ 4 & 1 & 8
\end{array} \right)
\xrightarrow{R_2-2R_1\to R_2}
\left( \begin{array}{ccc}
1 & 0 & 2\\ 0 & -1 & -1\\ 4 & 1 & 8
\end{array} \right)
\xrightarrow{R_3-4R_1\to R_3}
\left( \begin{array}{ccc}
1 & 0 & 2\\ 0 & -1 & -1\\ 0 & 1 & 0
\end{array} \right)
\xrightarrow{R_3+R_2\to R_3}
\left( \begin{array}{ccc}
1 & 0 & 2\\ 0 & -1 & -1\\ 0 & 0 & -1
\end{array} \right).
\]
\end{example}

\begin{definition}
Let $A$ be an $m\times n$ matrix over $\K$. The \emph{row rank} of $A$ is the dimension of the subspace of
$\K^n$ spanned by the rows of $A$. The \emph{column rank} of $A$ is the dimension of the subspace of $\K^m$
spanned by the columns of $A$.
\end{definition}

\begin{theorem}
For every matrix $A$ the row rank of $A$ equals the column rank of $A$.
\end{theorem}

\begin{definition}
The \emph{rank} of a matrix $A$, denoted $\rank(A)$, is its row rank (equivalently, its column rank).
\end{definition}

\begin{fact}
For any $m\times n$ matrix $A$ we have $\rank(A) = \rank(A^T)$ and $\rank(A)\leq \min(m,n)$: the rank is
bounded both by the number of rows and by the number of columns.
\end{fact}

\begin{example}
Consider $A=\left( \begin{smallmatrix} 1 & 2\\ 2 & 4 \end{smallmatrix} \right)$. Then
\[
\dim\big(\Span{\{(1,2),(2,4)\}}\big) = \dim\{a(1,2)+b(2,4)\mid a,b\in \R\}
= \dim\{(a+2b)(1,2)\mid a,b\in \R\} = \dim\big(\Span{\{(1,2)\}}\big)=1.
\]
Hence $\rank(A)=1$.
\end{example}

\begin{theorem}
If a matrix $A$ is row equivalent to a matrix $B$, then $\rank(A)=\rank(B)$.
\end{theorem}

\begin{theorem}
Let $A$ be a matrix. The number of nonzero rows in a row echelon form which is row equivalent to $A$ equals
$\rank(A)$.
\end{theorem}

\begin{example}
We compute the rank of a $3\times 4$ matrix by row reduction:
\[
A = \left( \begin{array}{cccc} 1 & 2 & 1 & 0 \\ 2 & 4 & 3 & 1 \\ 3 & 6 & 4 & 1 \end{array} \right)
\xrightarrow{\substack{R_2-2R_1\to R_2\\ R_3-3R_1\to R_3}}
\left( \begin{array}{cccc} 1 & 2 & 1 & 0 \\ 0 & 0 & 1 & 1 \\ 0 & 0 & 1 & 1 \end{array} \right)
\xrightarrow{R_3-R_2\to R_3}
\left( \begin{array}{cccc} 1 & 2 & 1 & 0 \\ 0 & 0 & 1 & 1 \\ 0 & 0 & 0 & 0 \end{array} \right).
\]
The row echelon form has two nonzero rows, so $\rank(A) = 2$.
\end{example}

\begin{theorem}
The dimension of the subspace $W$ of $\K^n$ spanned by vectors $\vv_1,\ldots, \vv_m$ equals the rank of the
matrix
\[
\left( \begin{array}{c} \vv_1\\ \vdots \\ \vv_m \end{array} \right).
\]
\end{theorem}

\begin{example}
What is the dimension of $W = \Span{\{(1,2,1),(2,4,3),(3,6,4)\}}\subseteq \R^3$? Writing the vectors as the
rows of a matrix and row reducing,
\[
\left( \begin{array}{ccc} 1 & 2 & 1 \\ 2 & 4 & 3 \\ 3 & 6 & 4 \end{array} \right)
\longrightarrow
\left( \begin{array}{ccc} 1 & 2 & 1 \\ 0 & 0 & 1 \\ 0 & 0 & 0 \end{array} \right),
\]
we get $\dim(W) = 2$. Moreover, the nonzero rows $(1,2,1)$ and $(0,0,1)$ of the echelon form constitute a
basis of $W$.
\end{example}

\section{Determinants}

To every square matrix $A \in \K^n_n$ over a field $\K$ one assigns a scalar from $\K$, denoted $|A|$ or
$\det(A)$ and called \emph{the determinant of $A$}. For matrices of sizes $1\times 1$, $2\times 2$ and
$3\times 3$ it is given by the explicit formulas
\[
\begin{aligned}
| a_{11} | &= a_{11},\\[2pt]
\left | \begin{array}{cc} a_{11} & a_{12}\\ a_{21} & a_{22} \end{array} \right | &= a_{11} a_{22} - a_{12} a_{21},\\[2pt]
\left | \begin{array}{ccc}
a_{11} & a_{12} & a_{13}\\ a_{21} & a_{22} & a_{23}\\ a_{31} & a_{32} & a_{33}
\end{array} \right | &=
a_{11} a_{22} a_{33} + a_{12} a_{23} a_{31}+a_{13} a_{21} a_{32}
- a_{13}a_{22}a_{31} -a_{12} a_{21} a_{33} - a_{11} a_{23} a_{32}.
\end{aligned}
\]

\begin{remark}
The $3\times 3$ formula is known as the \emph{rule of Sarrus}: copy the first two columns to the right of
the matrix; the three ``descending'' diagonal products are added and the three ``ascending'' ones
subtracted. Beware: the rule has \emph{no} analogue for $n\geq 4$ — larger determinants are computed by
the Laplace expansion defined below or, more efficiently, by row reduction.
\end{remark}

\begin{example}
\[
\left| \begin{array}{cc} 1 & 2\\ 3 & 4 \end{array} \right| = 1\cdot 4 - 2\cdot 3 = -2,
\]
and, applying the $3\times 3$ formula term by term,
\[
\left| \begin{array}{ccc} 2 & 1 & 3\\ 0 & 1 & -1\\ 1 & 2 & 1 \end{array} \right|
= 2\cdot 1\cdot 1 + 1\cdot(-1)\cdot 1 + 3\cdot 0\cdot 2
- 3\cdot 1\cdot 1 - 1\cdot 0\cdot 1 - 2\cdot(-1)\cdot 2
= 2 - 1 + 0 - 3 - 0 + 4 = 2.
\]
\end{example}

\begin{definition}[Laplace expansion]
Let $A$ be an $n\times n$ matrix and let $A_{ij}$ denote the matrix obtained from $A$ by deleting the $i$-th
row and the $j$-th column. For any fixed column index $j$ between $1$ and $n$,
\[
\det(A) = \sum_{i=1}^{n} (-1)^{i+j} a_{ij}\cdot \det(A_{ij}).
\]
This formula is the \emph{Laplace expansion along the $j$-th column}. In particular, for $j=1$,
\[
\det(A) = \sum_{i=1}^{n} (-1)^{i+1} a_{i 1}\cdot \det(A_{i 1}).
\]
\end{definition}

\begin{remark}
It is a non-trivial fact, which we accept without proof, that the result does not depend on the chosen
column $j$ — so the recursive definition is consistent, and for sizes $2$ and $3$ it reproduces the
explicit formulas above. Moreover, since $\det(A) = \det(A^T)$ (see the properties below), the determinant
may equally well be expanded along any fixed \emph{row}:
$\det(A) = \sum_{j=1}^{n} (-1)^{i+j} a_{ij}\cdot\det(A_{ij})$ for any fixed row index $i$. In practice one
expands along the row or column containing the most zeros.
\end{remark}

\begin{example}
We compute the determinant by expanding along the first column:
\[
\left | \begin{array}{ccc}
1 & 2 & 3\\ 4 & 5 & 6\\ 7 & 8 & 9
\end{array} \right |
= 1\cdot\left | \begin{array}{cc} 5 & 6\\ 8 & 9 \end{array} \right |
- 4 \cdot \left | \begin{array}{cc} 2 & 3\\ 8 & 9 \end{array} \right |
+ 7\cdot \left | \begin{array}{cc} 2 & 3\\ 5 & 6 \end{array} \right | = 0.
\]
\end{example}

\begin{fact}[Properties of determinants]
Let $A$ be a square matrix of size $n\times n$. Then
\begin{enumerate}
\item $\det(A) = \det(A^T)$;
\item $\det(A) = 0$ if $A$ has a zero row (or column) or two identical rows (or columns);
\item $\det(A) = 0$ if and only if $\rank(A)<n$;
\item $\det(A) = -\det(B)$ if $B$ is obtained from $A$ by a single row switch ($R_i\leftrightarrow R_j$);
\item $\det(A) = \det(B)$ if $B$ is obtained from $A$ by a single row addition
($R_i+k\cdot R_j\rightarrow R_i$);
\item $\det(B) = k\cdot \det(A)$ if $B$ is obtained from $A$ by the row scaling $kR_i\rightarrow R_i$.
\end{enumerate}
The last three properties also hold for the corresponding elementary column operations.
\end{fact}

\begin{example}
Using row operations to reach a triangular form:
\[
\left | \begin{array}{ccc}
1 & 2 & 3\\ 4 & 5 & 6\\ 7 & 8 & 9
\end{array} \right |
\overset{R_2-4 R_1}{=}
\left | \begin{array}{ccc}
1 & 2 & 3\\ 0 & -3 & -6\\ 7 & 8 & 9
\end{array} \right |
\overset{R_3-7 R_1}{=}
\left | \begin{array}{ccc}
1 & 2 & 3\\ 0 & -3 & -6\\ 0 & -6 & -12
\end{array} \right |
= 2\cdot \left | \begin{array}{ccc}
1 & 2 & 3\\ 0 & -3 & -6\\ 0 & -3 & -6
\end{array} \right | = 2\cdot 0 = 0.
\]
\end{example}

\begin{fact}
The determinant of an upper (or lower) triangular matrix equals the product of the entries on its main
diagonal. In particular this applies to diagonal matrices, and $\det( I_{n\times n} ) = 1$.
\end{fact}

\begin{example}
\[
\left| \begin{array}{ccc} 2 & 7 & -1\\ 0 & 3 & 5\\ 0 & 0 & -4 \end{array} \right| = 2\cdot 3\cdot (-4) = -24.
\]
Combined with row reduction, this fact gives an efficient general method of computing determinants: reduce
the matrix to a triangular form, keeping track of how each operation changes the determinant, exactly as in
the example above.
\end{example}

\begin{example}[A $4\times 4$ determinant]
Row additions do not change the determinant, so
\[
\left| \begin{array}{cccc}
1 & 2 & 0 & 1\\ 2 & 5 & 2 & 3\\ 1 & 2 & 1 & 2\\ 3 & 6 & 0 & 4
\end{array} \right|
\overset{\substack{R_2-2R_1\\ R_3-R_1\\ R_4-3R_1}}{=}
\left| \begin{array}{cccc}
1 & 2 & 0 & 1\\ 0 & 1 & 2 & 1\\ 0 & 0 & 1 & 1\\ 0 & 0 & 0 & 1
\end{array} \right| = 1\cdot 1\cdot 1\cdot 1 = 1.
\]
For larger matrices this approach is dramatically faster than a direct Laplace expansion.
\end{example}

\begin{fact}[Multiplicativity]
Let $A$ and $B$ be square matrices of size $n\times n$. Then $\det(A\cdot B) = \det(A)\cdot \det(B)$.
\end{fact}

\begin{warning}
The determinant is \emph{not} additive: in general $\det(A+B)\neq \det(A)+\det(B)$. Note also that
$\det(k\cdot A) = k^n\det(A)$ for an $n\times n$ matrix $A$: scaling the whole matrix scales \emph{each} of
its $n$ rows.
\end{warning}

\begin{remark}[Geometric meaning]
For a real $2\times 2$ matrix $A$, the absolute value $|\det(A)|$ is the area of the parallelogram spanned
by the two rows (equivalently, the two columns) of $A$; the sign of the determinant records the orientation
of this pair of vectors. Similarly, for a real $3\times 3$ matrix, $|\det(A)|$ is the volume of the
parallelepiped spanned by the rows.
\end{remark}

\section{Matrix inversion}

\begin{definition}
Let $A$ be a square matrix of size $n\times n$. A matrix $B$ of the same size is called an \emph{inverse} of
$A$ if
\[
A\cdot B = I_{n\times n}, \qquad B\cdot A = I_{n\times n}.
\]
If $A$ has an inverse, it is unique; it is denoted by $A^{-1}$.
\end{definition}

\begin{example}
For $A = \left ( \begin{smallmatrix} 1 & 0\\ 0 & 2 \end{smallmatrix} \right )$ one checks that
$A^{-1} = \left ( \begin{smallmatrix} 1 & 0\\ 0 & \tfrac12 \end{smallmatrix} \right )$.
\end{example}

\begin{warning}
Not every matrix is invertible.
\end{warning}

\begin{theorem}
A square matrix $A$ is invertible if and only if $\det(A)\neq 0$. If $A$ is invertible, then
\[
\det(A^{-1}) = \frac{1}{\det(A)}.
\]
\end{theorem}

\begin{proof}
We prove the second statement. Since $\det(I) = 1$ and $\det(A\cdot B) = \det(A)\cdot \det(B)$, and since
$A\cdot A^{-1} = I$ for invertible $A$, we get
\[
1 = \det(I) = \det(A\cdot A^{-1}) = \det(A)\cdot \det(A^{-1}). \qedhere
\]
\end{proof}

\begin{theorem}[Properties of the inverse]
Let $A,B$ be invertible matrices of size $n\times n$ and let $k\neq 0$ be a scalar. Then $A^{-1}$,
$A\cdot B$, $A^T$ and $k\cdot A$ are all invertible, and
\[
(A^{-1})^{-1} = A, \qquad (A\cdot B)^{-1} = B^{-1}\cdot A^{-1}, \qquad (A^T)^{-1} = (A^{-1})^T,
\qquad (k\cdot A)^{-1} = \tfrac1k\cdot A^{-1}.
\]
\end{theorem}

\begin{warning}
Note the reversed order in $(A\cdot B)^{-1} = B^{-1}\cdot A^{-1}$ — the same reversal as in the transpose
formula $(A\cdot B)^T = B^T\cdot A^T$.
\end{warning}

\begin{example}[The general linear group]
Recall the notion of a group from Chapter~\ref{ch:groups-fields}. By the theorem above, a product of two
invertible $n\times n$ matrices is invertible, and so is the inverse of an invertible matrix; moreover,
matrix multiplication is associative and $I$ is invertible. Hence the set of all invertible $n\times n$
matrices over $\K$, with matrix multiplication, is a group, called the \emph{general linear group}
$GL_n(\K)$. For $n\geq 2$ it is \emph{not} abelian; for instance, in $GL_2(\R)$,
\[
\begin{pmatrix} 1 & 1\\ 0 & 1\end{pmatrix}
\begin{pmatrix} 1 & 0\\ 1 & 1\end{pmatrix}
=\begin{pmatrix} 2 & 1\\ 1 & 1\end{pmatrix}
\neq
\begin{pmatrix} 1 & 1\\ 1 & 2\end{pmatrix}
=\begin{pmatrix} 1 & 0\\ 1 & 1\end{pmatrix}
\begin{pmatrix} 1 & 1\\ 0 & 1\end{pmatrix}.
\]
This is our first example of a non-abelian group.
\end{example}

\begin{fact}[Inversion by row reduction]
Let $A$ be a square matrix of size $n\times n$. Form the $n\times 2n$ matrix $C=(A\mid I_{n\times n})$ and row reduce
it to the form $(I_{n\times n}\mid B)$. Such a reduction is possible if and only if $A$ is invertible, and
the resulting matrix $B$ is the inverse of $A$.
\end{fact}

\begin{example}
\[
\left ( \begin{array}{cc|cc}
1 & 0 & 1 & 0 \\ 2 & 2 & 0 & 1
\end{array} \right )
\xrightarrow{R_2-2R_1\to R_2}
\left ( \begin{array}{cc|cc}
1 & 0 & 1 & 0 \\ 0 & 2 & -2 & 1
\end{array} \right )
\xrightarrow{\frac{1}{2}R_2\to R_2}
\left ( \begin{array}{cc|cc}
1 & 0 & 1 & 0 \\ 0 & 1 & -1 & \tfrac{1}{2}
\end{array} \right ),
\]
so $\left ( \begin{smallmatrix} 1 & 0\\ 2 & 2 \end{smallmatrix} \right )^{-1}
= \left ( \begin{smallmatrix} 1 & 0\\ -1 & \tfrac{1}{2} \end{smallmatrix} \right )$.
\end{example}

\begin{example}[A $3\times 3$ inverse]
We invert $A = \left ( \begin{smallmatrix} 1 & 0 & 1\\ 1 & 1 & 0\\ 0 & 1 & 1 \end{smallmatrix} \right )$:
\[
\left ( \begin{array}{ccc|ccc}
1 & 0 & 1 & 1 & 0 & 0\\ 1 & 1 & 0 & 0 & 1 & 0\\ 0 & 1 & 1 & 0 & 0 & 1
\end{array} \right )
\xrightarrow{R_2-R_1\to R_2}
\left ( \begin{array}{ccc|ccc}
1 & 0 & 1 & 1 & 0 & 0\\ 0 & 1 & -1 & -1 & 1 & 0\\ 0 & 1 & 1 & 0 & 0 & 1
\end{array} \right )
\xrightarrow{R_3-R_2\to R_3}
\left ( \begin{array}{ccc|ccc}
1 & 0 & 1 & 1 & 0 & 0\\ 0 & 1 & -1 & -1 & 1 & 0\\ 0 & 0 & 2 & 1 & -1 & 1
\end{array} \right )
\]
\[
\xrightarrow{\frac12 R_3\to R_3}
\left ( \begin{array}{ccc|ccc}
1 & 0 & 1 & 1 & 0 & 0\\ 0 & 1 & -1 & -1 & 1 & 0\\ 0 & 0 & 1 & \tfrac12 & -\tfrac12 & \tfrac12
\end{array} \right )
\xrightarrow{\substack{R_2+R_3\to R_2\\ R_1-R_3\to R_1}}
\left ( \begin{array}{ccc|ccc}
1 & 0 & 0 & \tfrac12 & \tfrac12 & -\tfrac12\\
0 & 1 & 0 & -\tfrac12 & \tfrac12 & \tfrac12\\
0 & 0 & 1 & \tfrac12 & -\tfrac12 & \tfrac12
\end{array} \right ),
\]
so
\[
A^{-1} = \frac12 \left ( \begin{array}{ccc} 1 & 1 & -1\\ -1 & 1 & 1\\ 1 & -1 & 1 \end{array} \right ).
\]
One can (and always should) verify the result by checking that $A\cdot A^{-1} = I$.
\end{example}

\begin{fact}[Inversion by determinants]
Let $A$ be an invertible square matrix and let $A_{k,m}$ denote the matrix obtained from $A$ by deleting its
$k$-th row and $m$-th column. Then
\[
A^{-1} = \frac{1}{|A|}\cdot
\left( \begin{array}{cccc}
(-1)^{1+1} |A_{11}| & (-1)^{1+2} |A_{12}| & \ldots & (-1)^{1+n} |A_{1n}| \\
(-1)^{2+1} |A_{21}| & (-1)^{2+2} |A_{22}| & \ldots & (-1)^{2+n} |A_{2n}| \\
\vdots & \vdots & \ddots & \vdots \\
(-1)^{n+1} |A_{n1}| & (-1)^{n+2} |A_{n2}| & \ldots & (-1)^{n+n}|A_{nn}|
\end{array} \right)^T.
\]
\end{fact}

\begin{example}
\[
\left ( \begin{array}{cc} 1 & 0 \\ 2 & 2 \end{array} \right )^{-1}
= \frac{1}{2} \left ( \begin{array}{cc}
|(2)| & -|(2)| \\ -|(0)| & |(1)|
\end{array} \right )^T
= \frac{1}{2} \left ( \begin{array}{cc} 2 & 0 \\ -2 & 1 \end{array} \right )
= \left ( \begin{array}{cc} 1 & 0 \\ -1 & \tfrac12 \end{array} \right ).
\]
\end{example}

\begin{fact}[Inverse of a $2\times 2$ matrix]
Let $A= \left ( \begin{smallmatrix} a & b \\ c & d \end{smallmatrix} \right )$ with
$\det(A) = ad-bc \neq 0$. Then
\[
A^{-1} = \frac{1}{ad-bc}\cdot \left( \begin{array}{cc} d & -b \\ -c & a \end{array} \right ).
\]
\end{fact}
